\naslov{Moduli}

\begin{definicija}
  Naj bo $R$ kolobar in $M$ neprazna množica.
  \pojem{Levi $R$-modul} je $(M, +, \cdot)$, kjer sta $+ : M \times M \to M$ in
  $\cdot : R \times M \to M$.
  Pri tem zahtevamo
  \begin{itemize}
  \item $(M, +)$ je Abelova grupa
  \item $r(m_1 + m_2) = r m_1 + r m_2$
  \item $(r_1 + r_2) m = r_1 m + r_2 m$
  \item $(r_1 r_2) m = r_1 (r_2 m)$
  \item $1 \cdot m = m$
  \end{itemize}
\end{definicija}

Analogno lahko definiramo desne module.
Vsak levi $R$-modul je pravzaprav tudi desni $R^{\text{opp}}$ modul, kjer je
$R^{\text{opp}}$ kolobar z enakim seštevanjem kot v $R$ in množenjem v obratnem
vrstnem redu.

\vprasanje{Definiraj modul.}

\begin{primer}
  Vektorski prostor nad poljem $F$ je levi (in desni) $F$-modul.
\end{primer}

\begin{primer}
  Abelova grupa (tu pisana aditivno) je $\Z$-modul za množenje
  \[
	n \cdot g = \underbrace{g + g + \ldots + g}_{n}.
  \]
\end{primer}

\begin{primer}
  Naj bo $I$ levi ideal kolobarja $R$.
  Potem je $I$ levi $R$-modul.
\end{primer}

\begin{primer}
  Če je $V$ vektorski prostor nad poljem $F$ in $R$ množica endomorfizmov $V \to
  V$, je $V$ levi $R$-modul za množenje
  \[
	A \cdot v = A(v).
  \]
\end{primer}

\vprasanje{Naštej nekaj enostavnih primerov modulov.}

\begin{definicija}
  Naj bo $M$ $R$-modul.
  Množica $N \subseteq M$ je \pojem{podmodul} v $M$, če je $N$ za podedovano
  seštevanje in množenje s skalarjem tudi modul.
  Oznaka $N \le M$.
\end{definicija}

\begin{definicija}
  Naj bo $M$ $R$-modul in $X \subseteq M$ neprazna množica.
  Najmanjšemu podmodulu, ki vsebuje $X$, pravimo \pojem{podmodul, generiran z
	množico $X$}.
  Oznaka $\sk{X}$.
  Modul je \pojem{končno generiran}, če obstaja končna $X$, da je $M = \sk{X}$.
  Če je modul generiran z enim samim elementom, je \pojem{ciklični}.
\end{definicija}

\begin{definicija}
  Naj bosta $N, M$ $R$-modula.
  Preslikava $\varphi : M \to N$ je \pojem{homomorfizem} $R$-modulov, če velja
  $\varphi(m_1 + m_2) = \varphi(m_1) + \varphi(m_2)$ in $\varphi(r m) = r
  \varphi(m)$.
  \pojem{Jedro} homomorfizma je množica
  \[
	\jedro \varphi = \{m \in M \such \varphi(m) = 0\},
  \]
  \pojem{slika} pa
  \[
	\im \varphi = \{\varphi(m) \such m \in M\}.
  \]
\end{definicija}

\begin{definicija}
  Naj bo $M$ $R$-modul in $N \le M$.
  Za relacijo
  \[
	m_1 \sim m_2 \nt m_1 - m_2 \in N
  \]
  definiramo \pojem{kvocientni modul} $M/_\sim = M/_N$.
\end{definicija}

\begin{izrek}
  Za kvocientne module veljajo naslednje lastnosti:
  \begin{itemize}
  \item če je $\varphi : M \to N$ homomorfizem, potem je $M/_{\jedro \varphi}
	\cong \im \varphi$,
  \item če sta $N, K$ podmodula v $M$, potem je $(N+K)/_K \cong N/_{N \presek
	  K}$,
  \item velja
	\[
	  \frac{M/_N}{L/_N} \cong M/_L.
	\]
  \end{itemize}
\end{izrek}

\begin{definicija}
  Naj bodo $N_1, \ldots, N_k$ $R$-moduli.
  \pojem{Zunanji direktni produkt} je $R$-modul $N_1 \times N_2 \times \cdots
  \times N_k$ s seštevanjem in množenjem s skalarjem po komponentah.
  Oznaka $N_1 \oplus N_2 \oplus \cdots \oplus N_k$.
\end{definicija}

\begin{opomba}
  Če imamo neskončno družino modulov, lahko naredimo podobno konstrukcijo, ki ji
  pravimo \pojem{kartezični produkt}.
\end{opomba}

\begin{definicija}
  Naj bodo $N_1, \ldots, N_k \le M$.
  Vsoti $N_1 + \cdots + N_k$ pravimo \pojem{notranja direktna vsota}, če za vsak
  $i$ velja $N_i \presek (N_1 + \cdots + N_{i-1} + N_{i+1} + \cdots + N_k) =
  \{0\}$.
\end{definicija}

\begin{opomba}
  Zunanja in notranja definicija sta ekvivalentni.
\end{opomba}

\begin{definicija}
  Naj bo $M$ $R$-modul in $X \subseteq M$.
  Pravimo, da je $X$ \pojem{baza} modula $M$, če je $M = \sk{X}$ in če za vsak
  $k \in \N$ in vse $x_1, \ldots, x_k \in X$ iz enakosti $r_1 x_1 + \cdots + r_n
  x_n = 0$ sledi $r_i = 0$.
\end{definicija}

\begin{primer}
  Obstajajo moduli, ki nimajo baz.
  Modul $M = \Z_n = \Z/_{n \Z}$ je $\Z$-modul.
  Recimo, da je nek $x + n\Z$ v bazi.
  Ampak $n(x + n \Z) = 0$, $n$ pa v $\Z$ ni enak $0$.
\end{primer}

\begin{definicija}
  Modul je \pojem{prost}, če ima bazo.
\end{definicija}

\vprasanje{Kaj je baza modula? Podaj primer modula, ki ni prost.}

\begin{izrek}
  Naj bo $M$ $R$-modul.
  Naslednje trditve so ekvivalentne:
  \begin{itemize}
  \item $M$ je prost $R$-modul.
  \item Obstaja indeksna množica $\Lambda$, da je $M$ kot $R$-modul izomorfen
	zunanji direktni vsoti kopij $R$-modula $R$:
	\[
	  M \cong \bigoplus_{\lambda \in \Lambda} R.
	\]
  \item Obstaja indeksna množica $\Lambda$ in podmoduli $M_\lambda \le M$,
	$M_\lambda \cong R$, da je
	\[
	  M = \bigoplus_{\lambda \in \Lambda} M_\lambda.
	\]
  \end{itemize}
\end{izrek}

\begin{proof}
  Druga in tretja točka sta očitno ekvivalentni.
  Naj bo $X$ baza modula $M$.
  Trdimo $M = \bigoplus_{x \in X} \sk{x}$.
  Ker je $M = \sk{X}$, je $M$ vsota podmodulov $\sk{x}$.
  Vzemimo $z \in \sk{x} \presek \sum_{y \ne x} \sk{y}$.
  Velja $z = rx = r_1 y_1 + \cdots + r_k y_k$, iz definicije baze pa sledi $r =
  r_i = 0$.
  Preveriti moramo še, da je $\sk{x} \cong R$.
  Ustrezen izomorfizem $R \to \sk{x}$ bo $\varphi(r) = rx$.

  Sedaj dokažimo, da iz druge točke sledi prva.
  Imamo izomorfizem
  \[
	\varphi : \bigoplus_{\lambda \in \Lambda} R \to M.
  \]
  Izberimo $e_\lambda$ kot $\Lambda$-terico, ki ima na mestu $\lambda$ enico,
  drugje pa $0$, in definirajmo $x_\lambda = \varphi(e_{\lambda})$.
  Preverimo lahko, da je $X = \{x_\lambda \such \lambda \in \Lambda\}$ baza.
\end{proof}

\vprasanje{Karakteriziraj prostost modula in dokaži karakterizacijo.}

\begin{posledica}
  Vsak $R$-modul je kvocient nekega prostega $R$-modula.
\end{posledica}

\begin{proof}
  Naj bo $M$ $R$-modul.
  Definiramo
  \[
	\varphi: \oplus_{m \in M} R \to M
  \]
  tako, da za element $e_m$ direktne vsote, ki ima na $m$-tem mestu enico,
  drugje pa nič, predpišemo $\varphi(e_m) = m$.
  Ti elementi tvorijo bazo, torej je s tem natanko določen epimorfizem
  $R$-modulov.
\end{proof}

\vprasanje{Pokaži, da je vsak modul kvocient prostega modula.}

\begin{primer}
  Podmodul prostega modula ni nujno prost.
  $\Z_8$ je prost $\Z_8$-modul, njegov podmodul $2 \Z_8 = \{0, 2, 4, 6\}$ pa ni
  prost.
\end{primer}

\vprasanje{Podaj primer podmodula prostega modula, ki ni sam prost.}

\begin{izrek}[univerzalna lastnost]
  $R$-modul $M$ je prost natanko tedaj, ko obstaja množica $X$ in preslikava
  $\iota: X \to M$, da velja: za vsak $R$-modul $N$ in vsako preslikavo $\kappa:
  X \to N$ obstaja natanko en homomorfizem $R$-modulov $f: M \to N$, da velja $f
  \circ \iota = \kappa$.
\end{izrek}

\begin{proof}
  V desno: $X$ naj bo baza $M$, $\iota$ pa vložitev, torej
  \begin{gather*}
	M = \oplus_{x \in X} R, \\
	\iota(x) = e_x.
  \end{gather*}
  Naj bo $N$ poljuben $R$-modul in $\kappa : X \to N$ poljubna preslikava.
  Recimo, da obstaja $f$, kot zahtevamo.
  Potem mora veljati $f \circ \iota = \kappa$, torej $f(e_x) = \kappa(x)$.
  Ker je ${\{e_x\}}_x$ baza, to enolično določi $f$, če le obstaja.
  Definiramo ga kot
  \[
	f\!\left( \sum_{x \in X} a_x e_x \right) = \sum_{x \in X} a_x \kappa(x),
  \]
  kar ustreza zahtevi.

  V levo:
  Dokazujemo, da je
  \[
	M \cong \oplus_{x \in X} R.
  \]
  Definiramo $\kappa$, kar nam po univerzalni lastnosti poda homomorfizem $f$,
  da diagram
  \[\begin{tikzcd}
	  X & {\oplus_x R} \\
	  M
	  \arrow["\kappa", from=1-1, to=1-2]
	  \arrow["\iota"', from=1-1, to=2-1]
	  \arrow["g"{description}, curve={height=6pt}, dotted, from=1-2, to=2-1]
	  \arrow["f"{description}, curve={height=6pt}, dotted, from=2-1, to=1-2]
	\end{tikzcd}\]
  komutira.
  Pokazali smo, da prosti moduli imajo univerzalno lastnost, torej obstaja
  natanko en homomorfizem $R$-modulov $g$ kot zgoraj, torej $g \circ \kappa =
  \iota$.
  Pokazati želimo, da sta $f$ in $g$ inverzni.
  Velja $f \circ g \circ \kappa = f \circ \iota = \kappa = \id \circ \kappa$.
  Oglejmo si torej diagram
  \[\begin{tikzcd}
	  X & {\oplus_x R} \\
	  {\oplus_x R}
	  \arrow["\kappa", from=1-1, to=1-2]
	  \arrow["\kappa"', from=1-1, to=2-1]
	  \arrow["{f \circ g}"{description}, curve={height=6pt}, from=1-2, to=2-1]
	  \arrow["{\id}"{description}, curve={height=6pt}, dotted, from=2-1, to=1-2]
	\end{tikzcd}\]
  Za oba homomorfizma diagram komutira, torej velja $f \circ g = \id$.
  Podobno v drugo smer.
\end{proof}

\vprasanje{Povej in dokaži izrek o univerzalni lastnosti prostih modulov.}

\begin{primer}
  Naj bo $F$ polje in $R$ množica endomorfizmov
  \[
	R = \End\!\left( \oplus_{n \in \N} F \right),
  \]
  torej neskončnih kvadratnih matrik elementov iz $F$.
  Kot $R$-modul je $R$ prost z bazo $\{\id\}$.
  Imamo pa tudi bazo iz dveh elementov:
  \begin{align*}
	B_1 =
	\begin{bmatrix}
	  1 & 0 & 0 & 0 & 0 & \cdots \\
	  0 & 0 & 1 & 0 & 0 & \cdots \\
	  0 & 0 & 0 & 0 & 0 & \cdots \\
	  \vdots & \vdots & \vdots & \vdots & \vdots & \ddots
	\end{bmatrix}
	&& B_2 =
	   \begin{bmatrix}
		 0 & 1 & 0 & 0 & 0 & \cdots \\
		 0 & 0 & 0 & 1 & 0 & \cdots \\
		 0 & 0 & 0 & 0 & 0 & \cdots \\
		 \vdots & \vdots & \vdots & \vdots & \vdots & \ddots
	   \end{bmatrix}
  \end{align*}
  Vsaka matrika $A$ je $R$-linearna kombinacija
  \[
	A = [a_1, a_3, a_5, \ldots] B_1 + [a_2, a_4, a_6, \ldots] B_2,
  \]
  torej je $\{B_1, B_2\}$ res baza (preverimo lahko, da sta res linearno
  neodvisni).
\end{primer}

\begin{definicija}
  Kolobar $R$ ima \pojem{lastnost enoličnega ranga}, če imajo vse baze
  poljubnega prostega $R$-modula isto kardinalnost.
\end{definicija}

\vprasanje{Povej primer kolobarja, ki nima lastnosti enoličnega ranga.}

\begin{primer}
  Polja in obsegi imajo lastnost enoličnega ranga.
\end{primer}

\begin{definicija}
  Naj bo $R$ kolobar, $I \edinka R$ in $M$ $R$-modul.
  Potem je
  \[
	IM = \{ \sum_{i=1}^k u_i m_i \such u_i \in I, m_i \in M, k \in \N \}
  \]
  podmodul v $M$.
\end{definicija}

\begin{lema}
  Če je $M$ prost $R$-modul z bazo $X$ in $I \triangleleft R$ različen od $R$,
  potem je $M/IM$ prost $R/I$-modul z bazo $X + IM$.
\end{lema}

\begin{proof}
  Očitno je $\sk{X + IM} = M/IM$, ker je $\sk{X} = M$.
  Linearna neodvisnost:
  \begin{align*}
	& \sum_i (r_i + I) (x_i + IM) = 0 + IM \\
	\nt& \sum_i (r_i x_i + IM) = 0 + IM \\
	\nt& \sum_i r_i x_i \in IM.
  \end{align*}
  Velja
  \[
	\sum_i r_i x_i = \sum_j u_j m_j = \sum_j u_j \left( \sum_k s_k x_k \right)
	= \sum_{j,k} \underbrace{u_j s_k}_{\in I} x_k.
  \]
  To sta dva razvoja po bazi $X$, torej $r_i = \sum_j u_j s_i \in I$, in $r_i +
  I = 0 + I$.
\end{proof}

\vprasanje{Dokaži: če je $M$ prost $R$-modul z idealom $I$, je $M/IM$ prost
  $R/I$-modul.}

\begin{lema}
  Ob oznakah prejšnje leme je $\abs{X} = \abs{X + IM}$.
\end{lema}

\begin{proof}
  Recimo, da je $x + IM = x' + IM$.
  Potem je $x - x' \in IM$, torej
  \[
	x - x' = \sum_j u_j x_j.
  \]
  Če je kakšen $u_j = 1$, je $I = R$, kar po predpostavki ne velja.
  Torej so vsi različni od $1$, in zaradi enoličnosti razvoja po bazi velja $x =
  x'$.
\end{proof}

\begin{izrek}
  Veljata naslednji točki o lastnosti enoličnega ranga:
  \begin{itemize}
  \item Naj bo $R$ kolobar.
	Recimo, da obstaja pravi ideal $I \edinka R$, da ima $R/I$ lastnost
	enoličnega ranga.
	Potem ima $R$ lastnost enoličnega ranga.
  \item Vsak komutativen kolobar ima lastnost enoličnega ranga.
  \end{itemize}
\end{izrek}

\begin{proof}
  Prva točka:
  Naj bo $M$ poljuben prost $R$-modul z bazama $X$ in $Y$.
  Vemo, da je $M/IM$ prost $R/I$-modul z bazama $X + IM$ in $Y + IM$.
  Ker ima $R/I$ lastnost enoličnega ranga, je $\abs{X + IM} = \abs{Y + IM}$,
  torej je po lemi $\abs{X} = \abs{Y}$.

  Druga točka:
  Naj bo $R$ komutativen kolobar.
  Če je $R$ polje, ima lastnost enoličnega ranga.
  Če pa ni polje, ima nek pravi ideal, ki je vsebovan v nekem maksimalnem idealu
  $I$.
  Potem je $R/I$ polje, torej ima lastnost enoličnega ranga.
  Po prvi točki ima tako $R$ lastnost enoličnega ranga.
\end{proof}

\vprasanje{Dokaži, da ima vsak komutativen kolobar lastnost enoličnega ranga.}

\podnaslov{Projektivni moduli}

\begin{definicija}
  $R$-modul $P$ je \pojem{projektiven}, če za vsak homomorfizem $R$-modulov $f:
  P \to M$ in vsak epimorfizem $R$-modulov $g: M' \to M$ obstaja homomorfizem
  $R$-modulov $h: P \to M'$, da naslednji diagram komutira:
  \[\begin{tikzcd}
	  P \\
	  M & {M'}
	  \arrow["f"', from=1-1, to=2-1]
	  \arrow["h", dotted, from=1-1, to=2-2]
	  \arrow["g", two heads, from=2-2, to=2-1]
	\end{tikzcd}\]
\end{definicija}

\vprasanje{Kaj je projektiven modul?}